\section{本章小结}

到本章为止，我们介绍完了概率论的基础知识。
我们从基础的概率问题出发，抽象到了纯数学的变量和函数，并讨论了两个变量相互作用的规律。
二维随机变量的分布情况会变得复杂，本章只是简单介绍，具体参见“教材\cite{book1,book2}”。

二维变量涉及的概念有：
\begin{itemize}
    \item 3个基本概念：联合分布，联合分布律，联合密度函数；
    \begin{align*}
    &F\left( x,y \right) =P\left\{ X\leqslant x,Y\leqslant y \right\} \\
    &P\left\{ X=x_i,Y=y_j \right\} =p_{ij} \\
    &f\left( x,y \right)
    \end{align*}
    \item 各自的边缘分布，这里要注意，边缘分布的计算过程丧失了其他变量的分布细节，所以已知边缘分布通常无法还原联合分布；
    \begin{align*}
    &\begin{cases}
	    F_X\left( x \right) =\underset{y\rightarrow +\infty}{\lim}F\left( x,y \right)\\
	    F_Y\left( y \right) =\underset{x\rightarrow +\infty}{\lim}F\left( x,y \right)\\
    \end{cases} \\
    &\begin{cases}
	    P\left\{ X=x_i \right\} =p_{i\cdot}=\sum_j{p_{ij}}\\
	    P\left\{ Y=y_j \right\} =p_{\cdot j}=\sum_i{p_{ij}}\\
    \end{cases} \\
    &\begin{cases}
    	f_X\left( x \right) =\int_{-\infty}^{+\infty}{f\left( x,y \right) \cdot dy}\\
    	f_Y\left( y \right) =\int_{-\infty}^{+\infty}{f\left( x,y \right) \cdot dx}\\
    \end{cases}
    \end{align*}
    \item 相互作用形成的条件分布，注意，一般不同条件下会获得不同的分布；
    \begin{align*}
    &\begin{cases}
	    F_{X|Y}\left( x \middle| y_0 \right) =P\left\{ X\leqslant x \middle| Y=y_0 \right\}\\
	    F_{Y|X}\left( y \middle| x_0 \right) =P\left\{ Y\leqslant y \middle| X=x_0 \right\}\\
    \end{cases} \\
    &\begin{cases}
    	P\left\{ X=x_i \middle| Y=y_{j_0} \right\} =\frac{p_{ij_0}}{p_{\cdot j_0}}\\
    	P\left\{ Y=y_j \middle| X=x_{i_0} \right\} =\frac{p_{i_0j}}{p_{i_0\cdot}}\\
    \end{cases} \\
    &\begin{cases}
    	f_{X|Y}\left( x \middle| y_0 \right) =\frac{f\left( x,y_0 \right)}{f_Y\left( y_0 \right)}\\
    	f_{Y|X}\left( y \middle| x_0 \right) =\frac{f\left( x_0,y \right)}{f_X\left( x_0 \right)}\\
    \end{cases}
    \end{align*}
    \item 相互的独立性。
    \begin{align*}
    &F\left( x,y \right) =F_X\left( x \right) \cdot F_Y\left( y \right)  \\
    &p_{ij}=p_{i\cdot}\cdot p_{\cdot j} \\
    &f\left( x,y \right) =f_X\left( x \right) \cdot f_Y\left( y \right)
    \end{align*}
\end{itemize}

关于二维变量引起的概念，可以和第一章中的条件概率、全概率、贝叶斯等概念结合，以前后呼应、融会贯通。
除此之外，还有二维随机变量的函数的分布和$n$维随机变量的分布，情况较为复杂，本笔记略。
最后谨记一点，我们所说的二维随机变量的概率，是两个事件都发生的概率。




